\documentclass[11pt,a4paper]{article}

\usepackage[margin=2cm]{geometry}
\usepackage[T1]{fontenc}
\usepackage{lmodern}
\usepackage{amsmath,amssymb}
\usepackage{graphicx}
\usepackage{booktabs}
\usepackage{multirow}
\usepackage{subcaption}
\usepackage{xcolor}
\usepackage{enumitem}
\usepackage[hidelinks]{hyperref}

\setlist{itemsep=1pt,topsep=3pt}
\graphicspath{{figures/}}
\let\oldus\_
\renewcommand{\_}{\oldus\allowbreak}
\emergencystretch=2em
\setlength{\textfloatsep}{8pt}
\setlength{\floatsep}{6pt}
\setlength{\abovecaptionskip}{3pt}

% Placeholder marker. Delete the \todo{} calls as sections get written.
\newcommand{\todo}[1]{\textcolor{red}{\textbf{[TODO:} #1\textbf{]}}}

\title{\vspace{-1.5cm}\textbf{Predicting Human Development Level from Economic Indicators}\\
\large DS203: Programming for Data Science, Group Project}
\author{\normalsize Vedant Iyengar (25B2141) \quad Rehan Mallik (25B0408)
  \quad Siddhartha Nair (25B3990)}
\date{\vspace{-0.6cm}}

\begin{document}
\maketitle

% =====================================================================
\section{Introduction}
% =====================================================================

The UNDP's Human Development Index (HDI) combines life expectancy, education
and income into one score, but it is published with a lag and is missing for
some countries: Monaco and North Korea have no 2023 value, and 22
territories tracked by the World Bank have none at all. World Bank
indicators on the economy, infrastructure and health are released more
often and for more economies. If they can predict a country's HDI tier
(three classes, Low / Medium-High / Very High), they can give an early
estimate where the HDI is late or absent, and show which measurable
characteristics are most closely associated with human development. We
compare a majority-class baseline, multinomial logistic regression, a
decision tree and a random forest, with and without PCA.

% =====================================================================
\section{Methodology}
% =====================================================================

\subsection{Dataset Download and Setup}

This process included us combining two different datasets, one from UNDP
HDI (composite-indices time series, 2023
HDI and region) and other from WBG API (\texttt{wbgapi}, the required
indicator list as per project doc, 2023 values). Merging
datasets by country (ISO-3 code). These are the countries/locations that had
to be dropped because of a lack of data: World Bank aggregates
(e.g.\ \texttt{WLD}, \texttt{HIC}) and 22 territories (e.g.\ ABW, BMU, PRI,
MAC, XKX). These 2 countries do not have a registered HDI\_2023 score: MCO
and PRK. Encoding our target response variable as per project requirement,
and creating \texttt{Services} and \texttt{Dominant\_Sector} as per project
doc, gives 195 countries. Source URLs are in \texttt{Data\_Dictionary.md}.

\subsection{Exploratory Data Analysis}
\label{sec:eda}

\paragraph{Dataset overview.} The dataset covers 195 countries across 14
columns, including 10 numerical predictors, 2 categorical predictors, 1
country identifier, and the target variable \texttt{Level}. The target
variable \texttt{Level} is not evenly distributed: 1.0 accounts for 47.69\%
(93 countries), 2.0 for 37.95\% (74), and 0.0 for 13.33\% (26).
\texttt{Dominant\_Sector} is derived from these shares, so including all of
them together could introduce redundancy.

\begin{table}[h]
\centering
\footnotesize
\caption{Features, missing counts (of 195) and class medians.}
\label{tab:features}
\begin{tabular}{@{}lrrrr@{}}
\toprule
Feature (unit) & Missing & \multicolumn{3}{c}{Median by class} \\
\cmidrule(l){3-5}
 & & 0 & 1 & 2 \\
\midrule
Internet users (\% pop.)        & 15 & 23.7 & 70.4 & 92.2 \\
Infant mortality (per 1{,}000)  & 2  & 45.7 & 17.7 & 4.0 \\
Access to electricity (\% pop.) & 0  & 45.4 & 99.3 & 100 \\
Agriculture (\% GDP)            & 12 & 29.0 & 9.9  & 2.1 \\
Services (\% GDP)               & 12 & 50.4 & 60.2 & 72.8 \\
Trade (\% GDP)                  & 28 & 55.3 & 76.2 & 107.3 \\
Industry (\% GDP)               & 11 & 22.2 & 26.8 & 23.7 \\
Inflation (annual \%)           & 20 & 7.1  & 6.0  & 5.0 \\
Unemployment (\% labour force)  & 18 & 4.5  & 5.3  & 5.2 \\
GDP growth (annual \%)          & 6  & 4.3  & 3.5  & 2.0 \\
Dominant\_Sector / Region       & 12 / 44 & \multicolumn{3}{c}{categorical} \\
\bottomrule
\end{tabular}
\end{table}

\paragraph{Missing value patterns.} Every labeled row with a missing
\texttt{Region} belongs to Class 2.0 (43 out of 43). This shows that
missingness in \texttt{Region} is systematic and strongly associated with
\texttt{Level = 2.0}, rather than occurring at random. Most macroeconomic
and digital-access metrics have noticeably higher missing rates in Class
0.0: \texttt{Inflation\_consumer\_prices},
\texttt{Individuals\_using\_the\_Internet}, and \texttt{Trade} are each
missing in 19.2\% of Class 0.0 countries (compared to 2.7\%, 0\%, and 9.5\%,
respectively, in Class 2.0). Imputing with a single global mean or median
could blur real differences between classes.

\paragraph{Numerical distributions \& outliers.}
\texttt{Inflation\_consumer\_prices} (17 upper outliers, up to 221.3\%) and
\texttt{Trade} (7 upper outliers, up to 402.6\%) have pronounced right
tails. \texttt{Access\_to\_electricity} piles up at the 100\% ceiling,
leaving 31 countries with lower electrification (down to 5.4\%) as
lower-tail outliers. Overall, several features exhibit strong skewness and
extreme values, which is worth accounting for when choosing scaling and
transformation steps for scale-sensitive models.

\paragraph{Patterns across HDI classes \& feature correlations.}
\texttt{Individuals\_using\_the\_Internet} and \texttt{Infant\_mortality}
show clear separation across classes. \texttt{Agriculture} decreases with
HDI, while \texttt{Services} increases (Table~\ref{tab:features}).
\texttt{Industry}, \texttt{Inflation\_consumer\_prices},
\texttt{Unemployment\_total}, and \texttt{GDP\_Growth} overlap considerably,
with differences mainly in their spread and outliers.
\texttt{Access\_to\_electricity} and \texttt{Individuals\_using\_the\_Internet}
are strongly positively correlated ($r=0.85$); both are negatively correlated
with \texttt{Infant\_mortality} ($r=-0.79$ and $-0.83$)
(Fig.~\ref{fig:eda}b). \texttt{Inflation\_consumer\_prices},
\texttt{GDP\_Growth}, and \texttt{Unemployment\_total} have virtually no
linear correlation with other variables (all $|r|<0.20$).

\paragraph{Categorical feature analysis.} EAP, LAC and SA are predominantly
Class 1, while ECA has a higher proportion of Class 2 countries. SSA has a
more mixed distribution, with a considerable share of Class 0 and Class 1
countries (Fig.~\ref{fig:eda}a). Services is the dominant sector for most
countries (173 out of 193 labelled observations). The smaller categories
have very few observations, so their distributions should be interpreted
carefully.

\begin{figure}[h]
  \centering
  \begin{subfigure}[b]{0.49\linewidth}
    \centering
    \includegraphics[width=\linewidth]{region_by_class.png}
    \caption{Class proportions by region.}
  \end{subfigure}\hfill
  \begin{subfigure}[b]{0.42\linewidth}
    \centering
    \includegraphics[width=\linewidth]{correlation_heatmap.png}
    \caption{Correlation of numerical features.}
  \end{subfigure}
  \caption{Exploratory analysis.}
  \label{fig:eda}
\end{figure}

\subsection{Preprocessing}

\begin{itemize}[leftmargin=*]
  \item \textbf{Split:} Removing rows where HDI is NaN. Dropping country
        name (it does not serve as a predictor). Splitting data into
        test/train/validation (stratified split i.e.\ roughly each class
        appears in the same amount in each split): 64/16/20, i.e.\
        123/31/39 countries, \texttt{random\_state=10}.
  \item \textbf{Numerical imputation:} Compared to mean imputation, median
        imputation is more resistant to outliers. As such we went with
        median imputation (our data has many outliers); we found medians
        w.r.t.\ train set. But, it is possible that a value being missing is
        important information to predict the HDI; to account for this
        possibility, we are adding missing indicators, i.e.\ for each
        numerical column, we are adding another column where each missing
        value gets a `1' and rest of the rows are `0'.
  \item \textbf{Categorical encoding:} missing \texttt{Region} and
        \texttt{Dominant\_Sector} are filled with ``Missing'', one-hot
        encoded, and reindexed to the training columns (dealing with unseen
        categories).
  \item \textbf{Scaling:} \texttt{RobustScaler} on numerical columns in a
        \texttt{ColumnTransformer} inside a \texttt{Pipeline}; it currently
        doesn't scale binary values.
  \item \textbf{PCA:} We compared the accuracy score of no PCA vs different
        number of components using PCA. To calculate the accuracy score we
        used LR and fit it on training set and predicted on the validation
        set.
\end{itemize}

\subsection{Models and Evaluation}

\section{Experiments}

\subsection{Preprocessing and PCA}

The class count percentages are approximately the same across train, test
and validation sets. To inspect principal component loadings, we have
plotted the first 2 components (that we got after PCA on the training set),
and their colour depends on their HDI (Fig.~\ref{fig:pca}a). Despite the
feature data being reduced to 2 features, we can still see that those with
similar colours are close together, especially the yellow cluster (although
there are outliers). Table~\ref{tab:loadings} gives the top features which
contributed the most to each component. N=8 is the ideal value of components
since it gave high accuracy at a lower value of N (Fig.~\ref{fig:pca}b).

\begin{figure}[h]
  \centering
  \begin{subfigure}[b]{0.4\linewidth}
    \centering
    \includegraphics[width=\linewidth]{pca_scatter.png}
    \caption{First two components (colour = class).}
  \end{subfigure}\hspace{1em}
  \begin{subfigure}[b]{0.4\linewidth}
    \centering
    \includegraphics[width=\linewidth]{pca_accuracy.png}
    \caption{Validation accuracy (LR) vs components.}
  \end{subfigure}
  \caption{PCA on the training set.}
  \label{fig:pca}
\end{figure}

\begin{table}[h]
\centering
\footnotesize
\caption{Top 3 features which contributed the most to each component
(by absolute loading).}
\label{tab:loadings}
\begin{tabular}{@{}llll@{}}
\toprule
PC1 & PC2 & PC3 & PC4 \\
\midrule
Inflation & Electricity & GDP growth & Unemployment \\
Trade & Internet & Unemployment & GDP growth \\
GDP growth & Infant mortality & Trade & Electricity \\
\bottomrule
\end{tabular}
\end{table}

\subsection{Model Comparison}

\subsection{Discussion and Conclusion}


\end{document}
